\documentclass[12pt]{article}
\usepackage[utf8]{inputenc}
\usepackage[T1]{fontenc}
\usepackage[margin=1in]{geometry}
\usepackage{setspace}
\usepackage{xcolor}
\usepackage[colorlinks=true, urlcolor=blue!60!black, linkcolor=blue!60!black]{hyperref}
\setlength{\parindent}{2em}
\setlength{\parskip}{1.25em}
\renewcommand{\baselinestretch}{1.0}
\title{\textit{Virginia State Board of Pharmacy v. Virginia Citizens Consumer Council}, 425 U.S. 748 (1976)}
\author{Case Brief}
\date{}
\begin{document}
\maketitle

\section{Case Metadata}
\begin{description}
  \item[Case Name:] \textit{Virginia State Board of Pharmacy v. Virginia Citizens Consumer Council, Inc.}
  \item[Citation:] 425 U.S. 748 (1976) (No.\ 74-895).
  \item[Argued / Decided:] Argued November 11, 1975; Decided May 24, 1976.
  \item[Vote:] 7--1 (Justice John Paul Stevens took no part in the consideration or decision of the case).
  \item[Author:] Justice Harry A.\ Blackmun (joined by Chief Justice Burger and Justices Brennan, Stewart, White, Marshall, and Powell).
  \item[Concurrences / Dissents:] Chief Justice Warren E.\ Burger and Justice Potter Stewart filed concurring opinions. Justice William H.\ Rehnquist filed a dissenting opinion.
  \item[Chapter / Topic:] Free Speech \& Press (Commercial Speech Doctrine; Listener's Right to Receive Information; Anti-Paternalism).
  \item[Primary Source URL:] \href{https://supreme.justia.com/cases/federal/us/425/748/}{Justia U.S. Supreme Court Center (425 U.S. 748)}.
\end{description}

\section{Facts and Procedural History}
Virginia Code Ann.\ \S\,54-524.35(3) (1974) declared a licensed pharmacist guilty of ``unprofessional conduct''---punishable by civil monetary penalty or license suspension or revocation---if the pharmacist ``publishes, advertises or promotes, directly or indirectly, in any manner whatsoever, any amount, price, fee, premium, discount, rebate or credit terms \dots\ for any drugs which may be dispensed only by prescription.'' Because only licensed pharmacists could dispense prescription drugs in Virginia, the statute effectively banned all prescription-drug price advertising, even though roughly 95\% of prescriptions were filled with prepackaged, manufacturer-prepared dosage forms rather than pharmacist-compounded mixtures.

In 1969, a three-judge federal court had upheld the same statute against a Due Process and Equal Protection challenge brought by a drug retailer and a pharmacist (\textit{Patterson Drug Co. v. Kingery}, 305 F.\ Supp.\ 821 (W.D. Va.\ 1969)). Here, however, the suit was brought not by pharmacists, but by prescription-drug \textit{consumers}: an individual Virginia resident whose medical conditions required daily prescription medication, along with two nonprofit organizations (the Virginia Citizens Consumer Council, Inc., and the Virginia State AFL-CIO). The parties stipulated that prescription drug prices varied dramatically within the same Virginia localities---for example, 40 Achromycin tablets in Richmond ranged from \$2.59 to \$6.00 (a 140\% spread), and tetracycline in Newport News--Hampton ranged from \$1.20 to \$9.00 (a 650\% spread)---and that, absent the ban, some Virginia pharmacists would advertise prices.

A three-judge U.S. District Court for the Eastern District of Virginia (373 F.\ Supp.\ 683 (1974)) distinguished \textit{Patterson Drug}, declared \S\,54-524.35(3) void under the First and Fourteenth Amendments, and enjoined its enforcement. The Virginia State Board of Pharmacy took a direct appeal to the Supreme Court.

\section{Issue(s) Presented}
\begin{enumerate}
  \item \textbf{Listener Standing:} Do consumers of prescription drugs have First Amendment standing to assert a reciprocal right to receive commercial price information from willing pharmacists?
  \item \textbf{Commercial Speech Protection:} Does purely commercial speech---expression that does ``no more than propose a commercial transaction''---fall wholly outside the protection of the First and Fourteenth Amendments, or does a State violate the First Amendment by banning truthful price advertising of lawful prescription drugs?
\end{enumerate}

\section{Holding}
Affirmed (7--1). First, where a willing speaker exists, First Amendment protection extends to the communication, its source, and its recipients alike; consumers therefore have a reciprocal First Amendment right to receive truthful advertising. Second, purely commercial speech is not wholly outside First Amendment protection, and a State may not completely suppress the dissemination of concededly truthful information about entirely lawful commercial activity out of fear of that information's effect on disseminators and recipients.

\section{Reasoning}
Writing for the Court, Justice Blackmun dismantled the categorical ``commercial speech exception'' that had persisted since \textit{Valentine v. Chrestensen}, 316 U.S.\ 52 (1942), and extended the trajectory begun in \textit{Bigelow v. Virginia}, 421 U.S.\ 809 (1975).

\begin{enumerate}
  \item \textbf{Reciprocal Right to Receive Information:} Relying on \textit{Lamont v. Postmaster General}, 381 U.S.\ 301 (1965), \textit{Kleindienst v. Mandel}, 408 U.S.\ 753 (1972), and \textit{Procunier v. Martinez}, 416 U.S.\ 396 (1974), the Court held that freedom of speech presupposes a willing speaker, and once a willing speaker exists (as stipulated here), the First Amendment protects both the speaker's right to communicate and the listener's reciprocal right to receive the message.

  \item \textbf{Value of Purely Commercial Speech:} Confronting speech that does ``no more than propose a commercial transaction'' (``I will sell you the X prescription drug at the Y price''), the Court rejected the premise that an advertiser's purely economic motive strips speech of First Amendment value. Drawing an analogy to labor-dispute speech (\textit{Thornhill v. Alabama}, 310 U.S.\ 88 (1940)), Blackmun emphasized three core interests: (a) the advertiser's economic interest; (b) the consumer's acute personal interest---especially for ``the poor, the sick, and particularly the aged,'' for whom price information can mean ``the alleviation of physical pain or the enjoyment of basic necessities'' (425 U.S.\ at 763--764); and (c) society's systemic interest in intelligent private resource allocation in a free-enterprise economy, which also informs democratic debate over how the economy should be regulated.

  \item \textbf{Rejection of State Paternalism:} The Pharmacy Board defended the ban as essential to preserving pharmacist professionalism, preventing destructive price-cutting, and protecting continuous pharmacist-patient monitoring. While acknowledging the State's strong interest in professional standards, the Court observed that Virginia already regulated pharmacy licensing and conduct directly. The advertising ban operated solely by keeping citizens ``in ignorance'' on the paternalistic assumption that consumers would unwisely choose low-cost, low-quality service. Rejecting this ``highly paternalistic approach,'' the Court held that the First Amendment itself chooses openness over enforced ignorance: Virginia may enforce professional standards directly or subsidize pharmacists, ``but it may not do so by keeping the public in ignorance of the entirely lawful terms that competing pharmacists are offering'' (425 U.S.\ at 770).

  \item \textbf{Preserved Regulatory Latitude (Footnote~24):} The Court stressed that commercial speech is not immune from all regulation. States remain free to enforce content-neutral time, place, and manner rules, prohibit false, deceptive, or misleading ads, and bar advertisements proposing unlawful transactions. In the famous Footnote~24 (425 U.S.\ at 771--772 n.\,24), Blackmun noted two ``common sense differences''---commercial speech's greater \textit{verifiability} (advertisers know their own products) and \textit{durability} (profit motive resists chilling)---that permit broader regulation of misleading claims, mandatory disclaimers, and potentially prior restraints, while Footnote~25 reserved judgment on advertising by lawyers and physicians.
\end{enumerate}

\section{Concurring and Dissenting Opinions}
\begin{itemize}
  \item \textbf{Chief Justice Burger, concurring (425 U.S.\ at 773--775):} Emphasized that 95\% of prescriptions involve prepackaged, standardized drugs, making pharmacists largely retail packagers in this context. He warned that advertising by the ``traditional learned professions of medicine or law,'' where individualized professional judgment predominates, poses far greater risks of deception and should not automatically follow from this holding.
  \item \textbf{Justice Stewart, concurring (425 U.S.\ at 775--781):} Wrote separately to reassure regulators that abandoning \textit{Chrestensen} does not jeopardize federal and state false-advertising laws. Unlike ideological speech---which requires ``breathing space'' for error under \textit{New York Times Co. v. Sullivan}---commercial price and product claims are empirically testable by the speaker, so banning false or deceptive ads promotes, rather than chills, accurate information.
  \item \textbf{Justice Rehnquist, dissenting (425 U.S.\ at 781--790):} Mounted a prescient, three-part critique. First, he contested consumer standing because Virginia did not bar consumers from inquiring by phone or publishing comparative price lists themselves. Second, he accused the majority of reviving \textit{Lochner}-era economic substantive due process under First Amendment guise, declaring that the Constitution does not require Virginia ``to hew to the teachings of Adam Smith'' (425 U.S.\ at 784). Third, he predicted that the Court's logic could not be confined to pharmacists---it would extend to lawyers, doctors, liquor, cigarettes, and aggressive direct-to-consumer drug promotion (offering hypothetical ads urging patients to ``Insist that your physician prescribe Demerol'' or ``Seconal''), overriding legitimate legislative judgments to discourage overconsumption of potentially harmful products.
\end{itemize}

\section{Key Quotations}
\begin{quote}
``If there is a right to advertise, there is a reciprocal right to receive the advertising, and it may be asserted by these appellees.''\\
--- Justice Blackmun, 425 U.S.\ at 757.
\end{quote}

\begin{quote}
``As to the particular consumer's interest in the free flow of commercial information, that interest may be as keen, if not keener by far, than his interest in the day's most urgent political debate. \dots\ So long as we preserve a predominantly free enterprise economy, the allocation of our resources in large measure will be made through numerous private economic decisions. It is a matter of public interest that those decisions, in the aggregate, be intelligent and well informed. To this end, the free flow of commercial information is indispensable.''\\
--- Justice Blackmun, 425 U.S.\ at 763, 765.
\end{quote}

\begin{quote}
``It is precisely this kind of choice, between the dangers of suppressing information, and the dangers of its misuse if it is freely available, that the First Amendment makes for us.''\\
--- Justice Blackmun, 425 U.S.\ at 770.
\end{quote}

\begin{quote}
``[T]here is certainly nothing in the United States Constitution which requires the Virginia Legislature to hew to the teachings of Adam Smith in its legislative decisions regulating the pharmacy profession.''\\
--- Justice Rehnquist, dissenting, 425 U.S.\ at 784.
\end{quote}

\section{Doctrinal Significance and Modern Legacy (Through 2026)}
\textit{Virginia State Board of Pharmacy} is the foundational decision of modern commercial speech jurisprudence, transforming advertising from an unprotected category into a central battleground of First Amendment law:
\begin{itemize}
  \item \textbf{Immediate Expansion to Professions:} Just as Rehnquist predicted, the Court swiftly invalidated blanket bans on truthful attorney price advertising in \textit{Bates v. State Bar of Arizona}, 433 U.S.\ 350 (1977), and contraceptive advertising in \textit{Carey v. Population Services Int'l}, 431 U.S.\ 678 (1977).
  \item \textbf{The \textit{Central Hudson} Framework:} Four years later, in \href{https://supreme.justia.com/cases/federal/us/447/557/}{\textit{Central Hudson Gas \& Electric Corp. v. Public Service Comm'n}}, 447 U.S.\ 557 (1980), the Court codified \textit{Virginia Pharmacy}'s principles into a four-part intermediate-scrutiny test: (1) the speech must concern lawful activity and not be misleading; (2) the asserted government interest must be substantial; (3) the regulation must directly advance that interest; and (4) it must not be more extensive than necessary.
  \item \textbf{Anti-Paternalism and Hardening Scrutiny:} In \textit{44 Liquormart, Inc. v. Rhode Island}, 517 U.S.\ 484 (1996), and \href{https://supreme.justia.com/cases/federal/us/533/525/}{\textit{Lorillard Tobacco Co. v. Reilly}}, 533 U.S.\ 525 (2001), the Court sharpened \textit{Virginia Pharmacy}'s anti-paternalism principle, striking down vice-advertising bans aimed at keeping consumers ignorant to reduce demand.
  \item \textbf{\textit{Sorrell v. IMS Health Inc.}, 564 U.S.\ 552 (2011)} (\href{https://supreme.justia.com/cases/federal/us/564/552/}{Justia}): Marking a major shift toward heightened scrutiny, the Court struck down a Vermont law barring pharmacies from selling prescriber-identifying data to pharmaceutical marketers. Writing for the Court, Justice Kennedy held that content- and speaker-based burdens on commercial expression trigger ``heightened judicial scrutiny''---and failed even \textit{Central Hudson}---explicitly invoking \textit{Virginia Pharmacy} to reject Vermont's effort to tilt the marketplace by silencing disfavored commercial speakers.
  \item \textbf{Modern Debates Through 2026 (``First Amendment Lochnerism'' and Compelled Disclosures):} Through 2026, \textit{Virginia Pharmacy}'s legacy sits at the center of two live controversies: (1) critics (echoing Justice Breyer's dissent in \textit{Sorrell} and Justice Kagan's warning about ``weaponizing the First Amendment'' in \textit{Janus v. AFSCME}, 585 U.S.\ 878 (2018)) argue that heightened scrutiny for commercial speech risks reviving Rehnquist's feared economic \textit{Lochnerism} against data-privacy, health, and financial regulations; and (2) courts actively police the boundary between permissible factual commercial disclosures (rooted in Footnote~24 and \textit{Zauderer v. Office of Disciplinary Counsel}, 471 U.S.\ 626 (1985)) and unconstitutional compelled speech (as in ongoing 2025--2026 litigation over mandatory climate, ESG, and platform disclosures).
\end{itemize}

\section{Pedagogical Critiques and Exam Hypotheticals}
\textbf{Pedagogical Critiques:}
\begin{itemize}
  \item \textit{Free-Market Constitutionalism vs.\ Democratic Self-Government:} Blackmun bridges commercial advertising to Alexander Meiklejohn's self-government theory by arguing that informed private economic choices inform public regulatory choices. Critics counter that this collapses the distinction between political participation and consumer preference, constitutionalizing laissez-faire market assumptions.
  \item \textit{Prescience of Rehnquist's Dissent:} Rehnquist's hypothetical ads for Demerol and Seconal anticipated the modern era of ubiquitous direct-to-consumer pharmaceutical advertising and the subsequent tension between First Amendment commercial speech rights and public-health efforts to curb prescription-drug overutilization.
\end{itemize}

\textbf{``Opposite Side'' Contrast (\textit{Virginia Pharmacy} vs.\ \textit{Williamson v. Lee Optical} and \textit{Ohralik}):}
Contrast \textit{Virginia Pharmacy} (1976) with \textit{Williamson v. Lee Optical Co.}, 348 U.S.\ 483 (1955), where the Court applied deferential rational-basis review under the Due Process Clause to uphold Oklahoma's ban on eyeglass-frame price advertising. Reframing an economic regulation as a restriction on \textit{truthful speech} flipped the judicial presumption from extreme deference to rigorous anti-paternalist scrutiny. Conversely, contrast \textit{Virginia Pharmacy} with \textit{Ohralik v. Ohio State Bar Ass'n}, 436 U.S.\ 447 (1978), where the Court upheld a ban on in-person lawyer solicitation of accident victims because face-to-face pressure---unlike impersonal price ads---carries acute risks of coercion, overreaching, and unverifiable deception.

\textbf{Exam Hypothetical:}
To combat rising rates of online sports-gambling addiction, a State where sports betting is legal enacts a statute prohibiting licensed sportsbooks from advertising ``promotional odds, signup bonuses, or payout percentages'' in any print or digital medium, while allowing sportsbooks to display those terms inside their mobile apps once a user creates an account. A consumer advocacy group representing recreational bettors and a licensed sportsbook sue under the First Amendment.
\begin{itemize}
  \item \textit{Challengers' Best Arguments:} Under \textit{Virginia Pharmacy}, the consumer group has standing via the reciprocal right to receive truthful commercial information. Because sports betting is lawful and the banned ads are truthful, the statute fails \textit{Central Hudson} and \textit{Sorrell}: it is a content- and speaker-based ban resting on forbidden paternalism (keeping bettors ignorant to dampen demand) rather than directly regulating betting limits or mandating addiction warnings.
  \item \textit{State's Best Arguments:} Distinguish \textit{Virginia Pharmacy} because the State is not banning \textit{all} access (terms remain available in-app, resembling a channel limit) and is targeting a vice activity with severe externalities rather than essential medicine. Under \textit{Central Hudson}, reducing gambling addiction is a substantial interest, promotional bonuses uniquely induce impulsive wagering (echoing \textit{Ohralik}'s concern with overreaching), and the restriction is tailored to promotional inducements rather than all speech about sportsbooks.
\end{itemize}

\end{document}
